\section{Vector Addition: \acr{AMD} \acr{ROCm}}

\textbf{\acr{HIP} equivalent (\acr{AMD}'s \acr{CUDA}-like \acr{API}):}
\begin{lstlisting}[style=ir]
__global__ void vectorAdd(float* a, float* b, float* c, int n) {
  int idx = hipBlockIdx_x * hipBlockDim_x + hipThreadIdx_x;
  if (idx < n) {
    c[idx] = a[idx] + b[idx];
  }
}
\end{lstlisting}

\textbf{TypeScript with AMDGPUCodegen:}
\begin{lstlisting}[style=ts]
import { AMDGPUCodegen } from "@euriklis/llvm-ir";

const codegen = new AMDGPUCodegen("vector_add");

// Same structure as CUDA, but use AMD intrinsics:
const workitemId = AMDGPUCodegen.rocm.workitemIdX("tid.x");
const workgroupId = AMDGPUCodegen.rocm.workgroupIdX("gid.x");
// Workgroup size is typically passed as kernel parameter

// Rest identical to CUDA version
\end{lstlisting}

\textbf{Key \acr{AMD} intrinsics:}
\begin{itemize}[noitemsep]
  \item \texttt{AMDGPUCodegen.rocm.workitemIdX/Y/Z}: Work-item ID (like \texttt{threadIdx})
  \item \texttt{AMDGPUCodegen.rocm.workgroupIdX/Y/Z}: Work-group ID (like \texttt{blockIdx})
  \item \texttt{AMDGPUCodegen.rocm.barrier()}: Work-group barrier
\end{itemize}

